% !TEX root = task-1-8-3.tex
% Editable task bank. Append new BankTask blocks at the end of the document.
\documentclass[12pt,a4paper]{article}
\usepackage{physics-textbook}
\usepackage{graphicx,needspace}
\usetikzlibrary{decorations.pathmorphing}
\newif\ifShowSolutions
\ifdefined\TutorEdition
  \ShowSolutionstrue
\else
  \ShowSolutionsfalse % Change to \ShowSolutionstrue to display solutions.
\fi
\newcounter{banktask}
\newsavebox{\BankTaskBox}
\newcommand{\BankTaskContent}[3]{%
  \noindent\textbf{\thebanktask.}\quad #2\par
  \ifShowSolutions
    \Needspace{50mm}% Keep the solution heading with its opening text or diagram.
    #3\par
    \ifdefined\TutorEdition
      \ifcsname TutorSource#1\endcsname
        \nopagebreak\smallskip{\small\csname TutorSource#1\endcsname\par}%
      \fi
    \fi
  \fi
}
\newcommand{\BankTask}[3]{%
  \refstepcounter{banktask}\par\medskip
  \sbox{\BankTaskBox}{\begin{minipage}{\linewidth}%
    \BankTaskContent{#1}{#2}{#3}\end{minipage}}%
  \ifdim\dimexpr\ht\BankTaskBox+\dp\BankTaskBox\relax>\dimexpr\textheight-20mm\relax
    \Needspace{35mm}\BankTaskContent{#1}{#2}{#3}%
  \else
    \Needspace{\dimexpr\ht\BankTaskBox+\dp\BankTaskBox\relax}%
    \noindent\usebox{\BankTaskBox}\par
  \fi
  \medskip
}
\begin{document}
\BookSubsection{1.8.3}{Третий закон Ньютона}


% Source sdamgia-219: https://phys-ege.sdamgia.ru/problem?id=219
\BankTask{sdamgia-219}{
Полосовой магнит массой $M$ поднесли к проводнику массой $m$, по которому течет ток. Сравните силу действия магнита на проводник $F_1$ с силой действия проводника на магнит $F_2.$

$1$) $F_1>F_2$

$2$) $F_1<F_2$

$3$) $F_1=F_2$

$4$) $\frac{F_1}{F_2} = \frac{M}{m}$
}{
\smallskip\noindent\textbf{Решение}\par
Согласно третьему закону Ньютона, сила $F_1,$ с которой магнит действует на проводник, равна по величине силе $F_2,$ с которой проводник действует на магнит.

Ответ: $3$.
}

% Source sdamgia-221: https://phys-ege.sdamgia.ru/problem?id=221
\BankTask{sdamgia-221}{
Солнце массой $M$ притягивает Землю массой $m$. Сравните силу действия Солнца на Землю $F_1$ с силой действия Земли на Солнце $F_2.$

$1$) $F_1>F_2$

$2$) $F_1<F_2$

$3$) $F_1=F_2$

$4$) $\frac{F_1}{F_2} = \frac{M}{m}$
}{
\smallskip\noindent\textbf{Решение}\par
Согласно третьему закону Ньютона, сила $F_1,$ с которой Солнце действует на Землю, равна по величине силе $F_2,$ с которой Земля действует на Солнце.

Ответ: $3$.
}

% Source sdamgia-223: https://phys-ege.sdamgia.ru/problem?id=223
\BankTask{sdamgia-223}{
Брусок массой $m$ лежит на доске массой $M$. Сравните силу действия доски на брусок $F_1$ с силой действия бруска на доску $F_2.$

$1$) $F_1>F_2$ 

$2$) $F_1<F_2$ 

$3$) $F_1=F_2$ 

$4$) $\frac{F_1}{F_2} = \frac{M}{m}$
}{
\smallskip\noindent\textbf{Решение}\par
Согласно третьему закону Ньютона, сила $F_1,$ с которой доска действует на брусок, равна по величине силе $F_2,$ с которой брусок действует на доску.

Ответ: $3$.
}

% Source sdamgia-242: https://phys-ege.sdamgia.ru/problem?id=242
\BankTask{sdamgia-242}{
Поршень массой $M$ удерживает в сосуде воздух массой $m$. Сравните силу действия поршня на воздух $F_1$ с силой действия воздуха на поршень $F_2.$

$1$) $F_1>F_2$ 

$2$) $F_1<F_2$ 

$3$) $F_1=F_2$ 

$4$) $\frac{F_1}{F_2} = \frac{M}{m}$
}{
\smallskip\noindent\textbf{Решение}\par
Согласно третьему закону Ньютона, сила $F_1,$ с которой поршень действует на воздух, равна по величине силе $F_2,$ с которой воздух действует на поршень.

Ответ: $3$.
}

% Source sdamgia-4480: https://phys-ege.sdamgia.ru/problem?id=4480
\BankTask{sdamgia-4480}{
Мальчик медленно поднимает гирю, действуя на нее с силой $100\,\text{Н}$. Гиря действует на руку мальчика с силой

$1$) меньше $100\,\text{Н}$, направленной вниз

$2$) больше $100\,\text{Н}$, направленной вниз

$3$) $100\,\text{Н}$, направленной вниз

$4$) $100\,\text{Н}$, направленной вверх
}{
\smallskip\noindent\textbf{Решение}\par
По третьему закону Ньютона, сила действия равна силе противодействия, при этом направлены эти силы противоположно. Следовательно, гиря действует на руку мальчика с силой $100\,\text{Н}$ и направлена вниз.

Ответ: $3$
}

% Source sdamgia-4515: https://phys-ege.sdamgia.ru/problem?id=4515
\BankTask{sdamgia-4515}{
Автомобиль массой $1000\,\text{кг}$ движется с постоянной по модулю скоростью по выпуклому мосту. Автомобиль действует на мост в верхней его точке с силой $F = 9000\,\text{Н}$. Сила, с которой мост действует на автомобиль, равна

$1$) $1000\,\text{Н}$ и направлена вертикально вверх

$2$) $19$ $000\,\text{Н}$ и направлена вертикально вниз

$3$) $9000\,\text{Н}$ и направлена вертикально вниз

$4$) $9000\,\text{Н}$ и направлена вертикально вверх
}{
\smallskip\noindent\textbf{Решение}\par
По третьему закону Ньютона, сила действия равна силе противодействия, при этом направлены эти силы противоположно. Следовательно, сила, с которой мост действует на автомобиль, равна $9000\,\text{Н}$ и направлена вертикально вверх.

Ответ: $4$.
}

% Source sdamgia-4725: https://phys-ege.sdamgia.ru/problem?id=4725
\BankTask{sdamgia-4725}{
 Мимо Земли летит астероид в направлении, показанном на рисунке пунктирной стрелкой. Вектор $\vec F_A$ показывает силу притяжения астероида Землей. Известно, что масса Земли в $10^5$ раз больше массы астероида. Вдоль какой стрелки ($1$ или $2$) направлена и чему равна по модулю сила, действующая на Землю со стороны астероида?

 

\begin{center}\begin{tikzpicture}[book diagram,>=stealth]
\draw[fill=gray!15](0,0)circle(.6);\fill(0,0)circle(1pt);
\draw[->](0,0)--(1,0)node[right]{$1$};\draw[->](0,0)--(.35,-1)node[right]{$2$};
\draw[fill=gray!25](2.05,0)--(2.2,.15)--(2.45,.1)--(2.4,-.13)--(2.2,-.1)--cycle;
\draw[->](2.2,0)--(1.45,0)node[midway,above]{$\vec F_A$};\draw[task projection,->](2.2,.1)--(1.85,1.1);
\end{tikzpicture}\end{center}
\noindent1) вдоль стрелки $1$, равна $F_A$\par
\noindent2) вдоль стрелки $1$, равна $10^5F_A$\par
\noindent3) вдоль стрелки $2$, равна $10^{-5}F_A$\par
\noindent4) вдоль стрелки $2$, равна $F_A$

}{
\smallskip\noindent\textbf{Решение}\par
Согласно третьему закону Ньютона, сила действия равна по величине силе противодействия, а направлены они противоположно. Поэтому сила, действующая на Землю со стороны астероида направлена вдоль стрелки $1$ и равна $F_{A}$

 

Ответ: $1$
}

% Source sdamgia-6915: https://phys-ege.sdamgia.ru/problem?id=6915
\BankTask{sdamgia-6915}{
Лошадь тянет телегу с постоянной скоростью. В соответствии с третьим законом Ньютона равны друг другу по модулю силы взаимодействия

$1$) лошади с телегой и телеги с лошадью

$2$) лошади с землей и телеги с землей

$3$) лошади с телегой и телеги с землей

$4$) лошади с телегой и лошади с землей
}{
\smallskip\noindent\textbf{Решение}\par
Третий закон Ньютона гласит, что тела действуют друг на друга с силами, равными по модулю и противоположными по направлению. Т. е. с какой силой лошадь тянет телегу, с такой же силой телега тянет лошадь. 

Ответ: $1$
}

% Source sdamgia-3784: https://phys-ege.sdamgia.ru/problem?id=3784
\BankTask{sdamgia-3784}{
Подъемный кран поднимает груз с постоянным ускорением. На груз со стороны каната действует сила, равная по величине $8\cdot10^3\,\text{Н}.$ На канат со стороны груза действует сила, которая

$1$) равна $8 \cdot  10 ^3\,\text{Н}$

$2$) меньше $8 \cdot  10 ^3\,\text{Н}$

$3$) больше $8 \cdot  10 ^3\,\text{Н}$

$4$) равна силе тяжести, действующей на груз
}{
\smallskip\noindent\textbf{Решение}\par
По третьему закону Ньютона, сила действия равна по величине силе противодействия. Поэтому сила, действующая со стороны груза на канат, равна по величине силе со стороны каната на груз, а именно, $8\cdot10^3\,\text{Н}.$

Ответ: $1$.
}

% Source sdamgia-3657: https://phys-ege.sdamgia.ru/problem?id=3657
\BankTask{sdamgia-3657}{
Два одинаковых груза массой $M=100\,\text{г}$ каждый подвешены на концах невесомой и нерастяжимой нити, перекинутой через невесомый блок с неподвижной осью. На один из них кладут перегрузок массой $m=20\,\text{г},$ после чего система приходит в движение. Найдите модуль силы $F$, действующей на ось блока во время движения грузов. Трением пренебречь.
}{
\smallskip\noindent\textbf{Решение}\par
\begin{center}\begin{tikzpicture}[book diagram,>=stealth]
\draw(-.5,2)--(.5,2);\foreach\x in{-.5,-.3,...,.5}\draw(\x,2)--++(.12,.12);
\draw(0,2)--(0,1);\draw(0,1)circle(.45);\fill(0,1)circle(1pt);
\draw(-.45,1)--(-.45,-.8)(.45,1)--(.45,-1.2);
\draw(-.58,-1.15)rectangle(-.32,-.8);\draw(.32,-1.55)rectangle(.58,-1.2);\draw(.3,-1.2)rectangle(.6,-1.05);
\draw[->](0,1)--(0,1.85)node[right]{$\vec F'$};\draw[->](0,1)--(0,.1)node[right]{$\vec F$};
\draw[->](-.45,1)--(-.45,.45);\draw[->](.45,1)--(.45,.45);
\draw[->](-.45,-.8)--(-.45,.15)node[left]{$\vec T_1$};\draw[->](.45,-1.05)--(.45,-.1)node[right]{$\vec T_2$};
\draw[->](-.45,-1.15)--(-.45,-1.9)node[left]{$M\vec g$};\draw[->](.45,-1.55)--(.45,-2.3)node[right]{$(M+m)\vec g$};
\draw[->](3.2,2)--(3.2,-2.2)node[below]{$X$};
\end{tikzpicture}\end{center}Запишем, на основании второго закона Ньютона, уравнения движения для обоих грузов в проекции на вертикальную ось, направленную вниз (см. рис.):

\[\begin{gathered}Ma_1=Mg - T_1,\\ ( M + m ) a_2= ( M + m ) g - T_2.\end{gathered}\]

Здесь через $a_1$ и $a_2$ обозначены проекции ускорений грузов $M$ и $( M + m )$ на вертикальную ось, а через $T_1$ и $T_2$ — проекции сил натяжения нити на ту же ось.

В силу условия задачи из-за нерастяжимости нити $a_1= - a_2= - a,$ а из-за невесомости блока и нити и отсутствия трения$T_1=T_2=T.$ Кроме того, в силу третьего закона Ньютона $F=F'=2T$ (здесь $F'$ — сила, действующая на блок со стороны его оси).

Из написанных уравнений получаем: 

\[\begin{gathered}a= \frac{m}{m + 2M} g,\\ F=2T=2Mg \frac{M + m}{M + ( \frac{m}{2} )} .\end{gathered}\]

Подставляя числовые данные и проверяя размерность, имеем: $F\approx 2{,}18\,\text{Н}.$

Ответ: $\displaystyle F=2Mg \frac{M + m}{M + ( \frac{m}{2} )} \approx 2{,}18\,\text{Н}.$
}

\end{document}
